% GENERATED FILE. Edit canonical lessons, then run npm run generate:books.
% canonical-source-hash: fnv1a64:b55b981dcaa3eb1f
% canonical-lessons: HI-C363-adhyayan, HI-C363-anusar, HI-C363-namuna, HI-C363-tulnakarna, HI-C363-akna

\chapter{A Study}
\label{ch:hi-a-study}
\hlchaptermodality{363}

\begin{chapteropening}
\textbf{By the end of this chapter:} \emph{I can report what a study found.}

Complete the last lesson of chapter 363: I can report what a study found.
\end{chapteropening}

\section[\texorpdfstring{\hi{अध्ययन} (adhyayan)}{अध्ययन (adhyayan)}]{\hi{अध्ययन} (adhyayan) — a study (a piece of research)}

\label{lesson:HI-C363-adhyayan}



\begin{quote}
Before the new one: say the Hindi for an analysis, then the Hindi for a trial.
\end{quote}

\subsection*{You'll want to know: \hi{अध्ययन}}
\textbf{\hi{अध्ययन}} — \emph{adhyayan} — \textquotedblleft{}a study (a piece of research)\textquotedblright{}.

\begin{quote}
\hi{मीरा} \hi{बच्चों} \hi{की} \hi{नींद} \hi{पर} \hi{अध्ययन} \hi{कर} \hi{रही} \hi{है।}
\end{quote}

\begin{quote}
\emph{Mīrā baccoṁ kī nīnd par adhyayan kar rahī hai.}
\end{quote}

\begin{quote}
Mira is doing a study on children's sleep.
\end{quote}

Masculine.

\begin{cousinweb}
\textbf{\hi{अध्ययन}} (\emph{adhyayan}) is a tatsama, borrowed from Sanskrit \emph{adhyayana}, \textquotedblleft{}reading, studying\textquotedblright{} (Monier-Williams).
\end{cousinweb}

\begin{grammarlens}[title={What you've built}]
One more word for giving reasons and opinions.
\end{grammarlens}

\subsection*{Your turn}
\begin{itemize}
\raggedright
  \item \emph{Say it:} \emph{adhyayan}
  \item \emph{Say it:} \emph{adhyayan}, once more
  \item \emph{Say it:} \emph{Mīrā baccoṁ kī nīnd par adhyayan kar rahī hai.} — Mira is doing a study on children's sleep.
  \item \emph{Recall:} say the Hindi for an analysis, then the Hindi for a trial, then say \emph{adhyayan} again
\end{itemize}

\begin{tcolorbox}[breakable,colback=teal!4,colframe=teal!35!black,title={Before you move on}]
What does \hi{अध्ययन} mean? (\textquotedblleft{}A study (a piece of research)\textquotedblright{}.) Say it once more.
\end{tcolorbox}
\section[\texorpdfstring{\hi{अनुसार} (anusār)}{अनुसार (anusār)}]{\hi{अनुसार} (anusār) — according to (\hi{के} \hi{अनुसार})}

\label{lesson:HI-C363-anusar}



\begin{quote}
Before the new one: say the Hindi for a trial, then the Hindi for a study.
\end{quote}

\subsection*{You'll want to know: \hi{अनुसार}}
\textbf{\hi{अनुसार}} — \emph{anusār} — \textquotedblleft{}according to (\hi{के} \hi{अनुसार})\textquotedblright{}.

\begin{quote}
\hi{नए} \hi{अध्ययन} \hi{के} \hi{अनुसार} \hi{चाय} \hi{सेहत} \hi{के} \hi{लिए} \hi{अच्छी} \hi{है।}
\end{quote}

\begin{quote}
\emph{nae adhyayan ke anusār cāy sehat ke liye acchī hai.}
\end{quote}

\begin{quote}
According to a new study, tea is good for your health.
\end{quote}

\begin{cousinweb}
\textbf{\hi{अनुसार}} (\emph{anusār}) is a tatsama, borrowed from Sanskrit \emph{anusāra}, \textquotedblleft{}following, conformity\textquotedblright{} (Monier-Williams).
\end{cousinweb}

\begin{grammarlens}[title={What you've built}]
One more word for giving reasons and opinions.
\end{grammarlens}

\subsection*{Your turn}
\begin{itemize}
\raggedright
  \item \emph{Say it:} \emph{anusār}
  \item \emph{Say it:} \emph{anusār}, once more
  \item \emph{Say it:} \emph{nae adhyayan ke anusār cāy sehat ke liye acchī hai.} — According to a new study, tea is good for your health.
  \item \emph{Recall:} say the Hindi for a trial, then the Hindi for a study, then say \emph{anusār} again
\end{itemize}

\begin{tcolorbox}[breakable,colback=teal!4,colframe=teal!35!black,title={Before you move on}]
What does \hi{अनुसार} mean? (\textquotedblleft{}According to (\hi{के} \hi{अनुसार})\textquotedblright{}.) Say it once more.
\end{tcolorbox}
\section[\texorpdfstring{\hi{नमूना} (namūnā)}{नमूना (namūnā)}]{\hi{नमूना} (namūnā) — a sample}

\label{lesson:HI-C363-namuna}



\begin{quote}
Before the new one: say the Hindi for a study, then the Hindi for according to.
\end{quote}

\subsection*{You'll want to know: \hi{नमूना}}
\textbf{\hi{नमूना}} — \emph{namūnā} — \textquotedblleft{}a sample\textquotedblright{}.

\begin{quote}
\hi{डाक्टर} \hi{ने} \hi{ख़ून} \hi{का} \hi{नमूना} \hi{लिया।}
\end{quote}

\begin{quote}
\emph{ḍākṭar ne khūn kā namūnā liyā.}
\end{quote}

\begin{quote}
The doctor took a blood sample.
\end{quote}

Masculine.

\begin{cousinweb}
\textbf{\hi{नमूना}} (\emph{namūnā}) is usually traced to Persian \emph{namūna}, \textquotedblleft{}a sample, a model\textquotedblright{}.
\end{cousinweb}

\begin{grammarlens}[title={What you've built}]
One more word for giving reasons and opinions.
\end{grammarlens}

\subsection*{Your turn}
\begin{itemize}
\raggedright
  \item \emph{Say it:} \emph{namūnā}
  \item \emph{Say it:} \emph{namūnā}, once more
  \item \emph{Say it:} \emph{ḍākṭar ne khūn kā namūnā liyā.} — The doctor took a blood sample.
  \item \emph{Recall:} say the Hindi for a study, then the Hindi for according to, then say \emph{namūnā} again
\end{itemize}

\begin{tcolorbox}[breakable,colback=teal!4,colframe=teal!35!black,title={Before you move on}]
What does \hi{नमूना} mean? (\textquotedblleft{}A sample\textquotedblright{}.) Say it once more.
\end{tcolorbox}
\section[\texorpdfstring{\hi{तुलना} \hi{करना} (tulnā karnā)}{तुलना करना (tulnā karnā)}]{\hi{तुलना} \hi{करना} (tulnā karnā) — to compare}

\label{lesson:HI-C363-tulnakarna}



\begin{quote}
Before the new one: say the Hindi for according to, then the Hindi for a sample.
\end{quote}

\subsection*{You'll want to know: \hi{तुलना} \hi{करना}}
\textbf{\hi{तुलना} \hi{करना}} — \emph{tulnā karnā} — \textquotedblleft{}to compare\textquotedblright{}.

\begin{quote}
\hi{दोनों} \hi{नमूनों} \hi{की} \hi{तुलना} \hi{करो।}
\end{quote}

\begin{quote}
\emph{donõ namūnoṁ kī tulnā karo.}
\end{quote}

\begin{quote}
Compare the two samples.
\end{quote}

Past, with \hi{ने} (he or she …): \hi{उसने} \hi{दोनों} \hi{की} \hi{तुलना} \hi{की}, \emph{usne donõ kī tulnā kī}.

\begin{cousinweb}
\textbf{\hi{तुलना} \hi{करना}} (\emph{tulnā karnā}) joins \hi{तुलना} to \hi{करना}, \textquotedblleft{}to do\textquotedblright{}. \hi{तुलना} is a tatsama, borrowed from Sanskrit \emph{tulanā}, \textquotedblleft{}equality, likeness\textquotedblright{} (Monier-Williams).
\end{cousinweb}

\begin{grammarlens}[title={What you've built}]
One more word for giving reasons and opinions.
\end{grammarlens}

\subsection*{Your turn}
\begin{itemize}
\raggedright
  \item \emph{Say it:} \emph{tulnā karnā}
  \item \emph{Say it:} \emph{tulnā karnā}, once more
  \item \emph{Say it:} \emph{donõ namūnoṁ kī tulnā karo.} — Compare the two samples.
  \item \emph{Recall:} say the Hindi for according to, then the Hindi for a sample, then say \emph{tulnā karnā} again
\end{itemize}

\begin{tcolorbox}[breakable,colback=teal!4,colframe=teal!35!black,title={Before you move on}]
What does \hi{तुलना} \hi{करना} mean? (\textquotedblleft{}To compare\textquotedblright{}.) Say it once more.
\end{tcolorbox}
\section[\texorpdfstring{\hi{आँकना} (ā̃knā)}{आँकना (ā̃knā)}]{\hi{आँकना} (ā̃knā) — to assess, to estimate (a value)}

\label{lesson:HI-C363-akna}



\begin{quote}
Before the new one: say the Hindi for a sample, then the Hindi for to compare.
\end{quote}

\subsection*{You'll want to know: \hi{आँकना}}
\textbf{\hi{आँकना}} — \emph{ā̃knā} — \textquotedblleft{}to assess, to estimate (a value)\textquotedblright{}.

\begin{quote}
\hi{बैंक} \hi{ने} \hi{घर} \hi{की} \hi{क़ीमत} \hi{आँकी।}
\end{quote}

\begin{quote}
\emph{baiṅk ne ghar kī qīmat ā̃kī.}
\end{quote}

\begin{quote}
The bank assessed the value of the house.
\end{quote}

Past, with \hi{ने} (he or she …): \hi{उसने} \hi{नुक़सान} \hi{आँका}, \emph{usne nuqsān ā̃kā}.

\begin{cousinweb}
\textbf{\hi{आँकना}} (\emph{ā̃knā}) is a tadbhava, inherited from Sanskrit \emph{aṅkayati}, \textquotedblleft{}marks\textquotedblright{} (CDIAL 104).
\end{cousinweb}

\begin{grammarlens}[title={What you've built}]
One more word for giving reasons and opinions.
\end{grammarlens}

\subsection*{Your turn}
\begin{itemize}
\raggedright
  \item \emph{Say it:} \emph{ā̃knā}
  \item \emph{Say it:} \emph{ā̃knā}, once more
  \item \emph{Say it:} \emph{baiṅk ne ghar kī qīmat ā̃kī.} — The bank assessed the value of the house.
  \item \emph{Recall:} say the Hindi for a sample, then the Hindi for to compare, then say \emph{ā̃knā} again
\end{itemize}

\begin{tcolorbox}[breakable,colback=teal!4,colframe=teal!35!black,title={Before you move on}]
What does \hi{आँकना} mean? (\textquotedblleft{}To assess, to estimate (a value)\textquotedblright{}.) Say it once more.
\end{tcolorbox}
